% !TEX root = ../main.tex
\chapter{电磁层析成像理论基础与传统成像方法}
\markboth{电磁层析成像理论基础与传统成像方法}{}

电磁层析成像是一种基于交变电磁场感应效应的非接触式无损检测技术。其基本思想是：通过传感器阵列向被测区域施加交变激励磁场，并测量被测介质对电磁场分布产生的扰动响应，进而反演敏感区域内部电导率、磁导率或与缺陷相关的等效电磁参数分布。对于高导电率、高磁导率金属构件而言，EMT 能够利用边界电磁响应感知表层及近表层伤损，因而在轨道交通金属部件检测中具有较好的应用潜力。

本章围绕后续研究所需的理论基础展开。首先介绍 EMT 系统的基本组成与工作流程；随后从麦克斯韦方程出发，分析 EMT 正问题的物理本质及其有限元求解方法；在此基础上，进一步讨论 EMT 逆问题的数学表达与病态性来源，并系统梳理基于灵敏度矩阵的传统成像方法及其局限性，为第三章所提出的 SCF-DeepONet 智能反演模型奠定理论基础。

\section{EMT 系统基本组成与工作流程}

一个典型的 EMT 系统通常由传感器阵列、激励与测量硬件平台以及上位机成像与分析软件三部分构成，其基本工作流程如图 \ref{fig:emt_basic_composition} 所示。

\begin{figure}[htbp]
    \centering
    \includegraphics[width=0.72\textwidth]{figures/EMT_system.jpg}
    \figcaption{EMT 系统基本组成}{Basic composition of an EMT system}
    \label{fig:emt_basic_composition}
\end{figure}

其中，传感器阵列负责与被测对象进行电磁耦合，既可产生一次激励磁场，也可感知目标区域对该磁场的扰动响应；激励与测量硬件平台负责完成激励信号产生、通道切换、微弱信号调理、锁相解调与模数转换等任务；上位机则负责对采集到的边界响应数据进行预处理、特征组织与成像求解，最终输出可视化检测结果。

需要指出的是，EMT 系统的具体硬件实现形式会因应用对象而有所不同。对于工业管道、封闭腔体等目标，常采用环形阵列；而对于钢轨、道岔等表面平整且空间受限的金属构件，则更适合采用平面阵列或柔性共形阵列，以尽可能减小提离距离并提高缺陷响应的可测性。本文面向道岔伤损检测所设计的系统属于后者，其具体硬件实现将在第五章中展开。

从测量机理上看，EMT 的一轮完整数据采集通常遵循“激励—响应—切换—重组”的流程。首先由激励模块向指定线圈施加交变电流，在线圈附近及被测金属中建立时变电磁场；随后，其余接收线圈测得由介质分布和伤损共同调制后的感应响应；通过多路选通网络依次切换激励与接收通道组合，便可获得一组多视角边界测量数据。后续成像任务正是基于这些有限边界观测值，反演敏感区域内部的电磁参数分布或与缺陷相关的等效表征。

\section{EMT 的数学与物理基础}

\subsection{从投影思想到 EMT 离散成像模型}

层析成像技术的共同目标，是通过边界或外部可测信息反推出目标内部的空间分布。对于 X 射线 CT 等“硬场”成像问题，探测射线可近似看作沿直线传播，因此可借助雷登变换建立从内部函数到外部投影的解析映射关系，并通过反投影实现图像重建\citep{radon20051}。EMT 虽然不满足直线传播条件，但其中“由外部响应反演内部结构”的思想，对其成像模型构建仍具有启发意义。

与 CT 不同，EMT 所依赖的低频交变电磁场在空间中呈弥散分布，具有明显软场特征。电磁场线会随介质电导率和磁导率分布变化而畸变，因此边界响应不再是简单的直线积分结果，而是电磁场传播、涡流耦合和材料参数分布共同作用的综合体现。也正因如此，EMT 难以获得类似 CT 的严格解析逆公式，通常需要将连续场问题离散化，并借助灵敏度矩阵建立近似线性模型。

设将敏感区域离散为 $N$ 个像素或网格单元，将单次测量得到的边界响应组织为 $M$ 维观测向量，则 EMT 的离散成像模型可写为
\begin{equation}
\label{eqn:EMT_linear}
\Delta V = S\Delta G + e ,
\end{equation}

式中，$\Delta V$ 表示相对于参考状态的边界响应变化量，$\Delta G$ 表示待重建区域内部电磁参数变化量，$S$ 为灵敏度矩阵，$e$ 为测量噪声与模型误差。该模型虽然是对真实非线性问题的近似，但构成了传统 EMT 成像算法的基本数学基础。

\subsection{EMT 正问题的电磁学描述}

EMT 的正问题是指：在已知几何结构、材料参数分布以及激励条件的前提下，求解场域内部电磁场分布，并进一步计算边界传感器上的响应信号。对于 EMT 而言，正问题是逆问题求解、灵敏度矩阵计算以及仿真数据生成的基础。

EMT 系统中的电磁场演化满足麦克斯韦方程组。其微分形式可表示为
\begin{align}
\label{eqn:maxwell_faraday}
\nabla \times \mathbf{E} &= -\frac{\partial \mathbf{B}}{\partial t} \\
\label{eqn:maxwell_ampere}
\nabla \times \mathbf{H} &= \mathbf{J} + \frac{\partial \mathbf{D}}{\partial t} \\
\label{eqn:maxwell_gauss_e}
\nabla \cdot \mathbf{D} &= \rho \\
\label{eqn:maxwell_gauss_b}
\nabla \cdot \mathbf{B} &= 0
\end{align}

其中，$\mathbf{E}$、$\mathbf{H}$、$\mathbf{D}$、$\mathbf{B}$ 和 $\mathbf{J}$ 分别表示电场强度、磁场强度、电位移、磁感应强度和电流密度。材料本构关系为
\begin{equation}
\label{eqn:constitutive}
\mathbf{D}=\varepsilon\mathbf{E}, \qquad
\mathbf{B}=\mu\mathbf{H}, \qquad
\mathbf{J}=\sigma\mathbf{E}
\end{equation}

在 EMT 常用工作频段内，电磁波波长通常远大于测量区域尺寸，同时在高导电率金属中，位移电流项相对于传导电流项通常较小，因此可采用低频准静态近似。在此条件下，EMT 场问题可通过磁矢势 $\mathbf{A}$ 与标量电位 $\phi$ 进行描述。由
\begin{equation}
\label{eqn:B_from_A}
\mathbf{B}=\nabla\times\mathbf{A}
\end{equation}

以及
\begin{equation}
\label{eqn:E_from_A_phi}
\mathbf{E}=-\nabla\phi-\frac{\partial \mathbf{A}}{\partial t}
\end{equation}

可得导体中的感应电流密度为
\begin{equation}
\label{eqn:eddy_current}
\mathbf{J}_e=\sigma\mathbf{E}
=-\sigma\nabla\phi-\mathrm{j}\omega\sigma\mathbf{A}
\end{equation}

进一步结合安培环路定律与适当规范条件，可得到描述 EMT 时谐场问题的控制方程。其一般形式可写为
\begin{equation}
\label{eqn:maxwell_control}
\nabla^2 \mathbf{A}
-\mathrm{j}\omega\mu\sigma\mathbf{A}
-\mu\sigma\nabla\phi
=
-\mu\mathbf{J}_s
\end{equation}

其中，$\mathbf{J}_s$ 为外部激励源电流密度。式 \eqref{eqn:maxwell_control} 表明，只要给定介质参数分布、激励条件及边界条件，理论上便可求得空间中的电磁场分布，并进一步计算接收线圈感应电压或等效阻抗变化。

\section{EMT 正问题的有限元求解}

\subsection{有限元建模思路}

尽管 EMT 正问题具有明确的偏微分方程描述，但由于实际系统中几何边界复杂、材料参数空间分布不均匀且耦合关系强烈，通常难以获得解析解。因此，工程中普遍采用有限元法（Finite Element Method, FEM）对 EMT 正问题进行数值求解。

有限元求解 EMT 正问题的基本流程包括几何建模、材料赋值、边界条件设置、网格划分、方程离散及数值求解。首先，根据传感器阵列和被测对象的真实几何尺寸建立仿真模型；随后为不同区域赋予相应的电导率和磁导率等材料参数；再对外部开放边界设置适当截断条件，将无限空间问题转化为有限计算域问题；最后通过网格划分和方程离散，求解场域内的磁矢势、电场、磁场以及感应电压等量。

\subsection{EMT 有限元仿真的关键问题}

在 EMT 数值计算中，网格质量和边界处理方式对求解精度具有重要影响。对于高导电率铁磁金属而言，由于存在明显趋肤效应，交变电磁场主要集中于材料表层，内部场量衰减较快。因此，在有限元仿真中通常需要对金属表面、传感器附近以及缺陷区域进行局部网格加密，而对远离目标的空气域采用相对稀疏的网格，以在计算精度和求解效率之间取得平衡。

此外，EMT 正问题的后处理结果不仅包括场分布云图，还包括线圈感应电压、互感变化量以及灵敏度矩阵计算所需的中间场量。这些结果一方面可用于分析不同缺陷形态对边界响应的影响规律，另一方面也为传统成像算法和后续深度学习模型的数据集构建提供了基础支撑。

\section{EMT 逆问题及其病态性}

\subsection{EMT 逆问题的数学表达}

与正问题相对应，EMT 的逆问题是指：在给定边界测量响应的条件下，反推出敏感区域内部电磁参数分布。设 EMT 正向映射记为非线性算子 $F$，则有
\begin{equation}
\label{eqn:forward_operator}
V = F(\sigma,\mu)
\end{equation}

逆问题即为求解 $F^{-1}$。然而，由于该映射通常高度非线性且难以解析求逆，工程上一般在参考状态附近对其进行局部线性化。令 $\sigma_0$ 和 $\mu_0$ 为参考材料分布，则在小扰动条件下有
\begin{equation}
\label{eqn:linearization}
\Delta V \approx F'(\sigma_0,\mu_0)\Delta G
\end{equation}

将连续场进一步离散化，便得到式 \eqref{eqn:EMT_linear} 所示的矩阵模型。求解 $\Delta G$，即构成传统 EMT 图像重建问题的基本形式。

\subsection{逆问题的病态性来源}

根据 Hadamard 关于适定问题的定义，一个良态问题应满足解存在、解唯一且对输入扰动稳定。然而，EMT 逆问题通常难以同时满足上述条件，是典型的病态逆问题。其原因主要体现在以下三个方面。

首先，EMT 属于软场问题。电磁场分布会随着介质参数变化而动态改变，因此边界响应与内部参数之间并非严格线性关系。式 \eqref{eqn:EMT_linear} 仅是局部线性近似，在高导电率、高磁导率铁磁材料背景下，这种近似误差往往更加明显。

其次，EMT 成像面临严重欠定性。实际系统中的独立测量通道数量有限，而为了获得较高空间分辨率，待重建区域往往需要划分为大量网格单元，因此未知量维数通常远大于观测量维数，导致同一组测量数据可能对应多种不同的内部参数分布。

最后，灵敏度矩阵通常具有较大的条件数，对测量噪声和模型误差极为敏感。特别是在目标区域深处或远离传感器的位置，灵敏度显著衰减，使得深层信息容易被噪声淹没。这也是传统方法在复杂金属场景中容易出现伪影、边界模糊和位置偏移的重要原因。

综上，强非线性、欠定性与噪声放大效应共同决定了 EMT 逆问题求解的困难程度。为了在此基础上获得稳定、可解释的重建结果，传统研究发展出了基于灵敏度矩阵的多种图像重建方法。

\section{基于灵敏度矩阵的传统成像方法}

\subsection{灵敏度矩阵的计算}

灵敏度矩阵是传统 EMT 成像方法的核心。其元素反映了某一网格单元电磁参数变化对某一测量通道响应的影响程度。常用灵敏度计算方法主要包括微小扰动法和场量提取法\citep{griffiths2001magnetic}。

微小扰动法基于灵敏度定义，通过在参考场域中逐个对网格单元施加微小已知扰动，并计算边界响应变化量来估计对应灵敏度系数。该方法物理意义直观，也可用于实验标定，但计算代价较高。

场量提取法则基于电磁互易原理，通过提取激励场和伴随场在局部区域内的分布关系，直接计算该区域的灵敏度值。该方法计算效率较高，已成为有限元仿真环境中构造灵敏度矩阵的常用途径。

基于电磁互易原理，有限元环境中灵敏度矩阵的元素可通过激励场与伴随场的内积直接计算，无需逐点施加扰动。具体地，第 $i$ 个测量通道对第 $j$ 个网格单元电磁参数变化的灵敏度可写为
\begin{equation}
\label{eqn:sensitivity_reciprocity}
S_{ij}
=
\frac{\partial V_i}{\partial G_j}
=
\int_{\Omega_j}
\bigl( \nabla \times \mathbf{A}^{(m)} \bigr)
\cdot
\bigl( \nabla \times \mathbf{A}^{(a,n)} \bigr)
\,\mathrm{d}\Omega
\end{equation}

其中 $\mathbf{A}^{(m)}$ 为第 $m$ 个激励线圈产生的激励场，$\mathbf{A}^{(a,n)}$ 为第 $n$ 个接收线圈对应的伴随场，$\Omega_j$ 为第 $j$ 个网格单元的积分区域，$V_i$ 为由激励--接收对 $(m,n)$ 构成的第 $i$ 个测量通道的响应。该公式表明，灵敏度矩阵的每一列均可通过一次正问题求解（伴随场）获得，相比微小扰动法具有显著的计算效率优势。

\subsection{线性反投影算法}

线性反投影（Linear Back Projection, LBP）是 EMT 中最基础的图像重建方法。该方法借鉴层析成像中“反投影”的思想，直接利用灵敏度矩阵的转置对边界测量数据进行反向映射。其表达式可写为
\begin{equation}
\label{eqn:LBP}
\Delta G =
\frac{S^T\Delta V}{\|S^T\Delta V\|}
\end{equation}

LBP 的优点是计算简单、速度快，易于实现实时成像；但其根本缺点在于并未真正求解逆问题，而只是以 $S^T$ 粗略代替逆映射，因此重建图像往往存在较强伪影和明显模糊，只能用于获得目标的大致位置分布。

\subsection{正则化方法}

为改善 EMT 逆问题中的噪声放大与解不稳定问题，通常在最小二乘框架中引入正则化项，对解施加先验约束。其中，典型方法为 Tikhonov 正则化，其目标函数写为
\begin{equation}
\label{eqn:Tikhonov}
J(\Delta G)
=
\min_{\Delta G}
\left\{
\|\Delta V-S\Delta G\|_2^2
+
\lambda\|R\Delta G\|_2^2
\right\}
\end{equation}

式中，第一项为数据保真项，第二项为正则化约束项，$\lambda$ 为正则化参数，$R$ 为正则化矩阵。其解析解为
\begin{equation}
\label{eqn:Tikhonov_solution}
\Delta G =
(S^TS+\lambda R^TR)^{-1}S^T\Delta V
\end{equation}

与 LBP 相比，正则化方法能够在一定程度上提升重建稳定性并抑制噪声，但当正则化程度过强时，又会引入过度平滑，导致边界细节和局部缺陷特征被削弱。

\subsection{迭代方法}

除直接解析法外，EMT 中还常采用 Landweber 迭代、共轭梯度法以及 Newton-Raphson 等迭代求解方法。以 Landweber 方法为例，其更新公式为
\begin{equation}
\label{eqn:Landweber}
\Delta G_{k+1}
=
\Delta G_k
-
\alpha_k S^T(S\Delta G_k-\Delta V)
\end{equation}

其中，$\alpha_k$ 为迭代步长。该类方法通过反复计算残差并更新解向量，能够在一定程度上提高重建精度，但通常计算开销较大，对初值、步长和停止准则较为敏感，因此在实时在线检测场景下应用受到一定限制。

\subsection{传统方法的局限性}

尽管基于灵敏度矩阵的方法在 EMT 研究中具有重要地位，但当其应用于高导电率、强磁性金属伤损检测时，会暴露出明显局限。

首先，传统方法普遍建立在线性微扰近似之上，而铁磁性金属中的电磁场分布具有更强非线性和更显著趋肤效应，导致线性模型误差增大。其次，单一灵敏度矩阵所包含的信息维度有限，难以充分表征复杂伤损对多激励电磁场的空间扰动。再次，传统方法普遍对噪声和模型误差较为敏感，尤其是在微弱信号和复杂工业干扰条件下，容易出现伪影增强、边界失真和重建不稳定等问题。

因此，传统 EMT 成像方法虽然具有明确物理背景和数学基础，但在复杂铁磁金属目标检测中的精度、鲁棒性和泛化能力仍然有限。如何摆脱对线性灵敏度模型的强依赖，并构建更适合复杂 EMT 逆问题的数据驱动反演模型，成为后续研究的关键方向。

\section{从传统方法到算子学习：深度学习范式的选择}

上述分析表明，传统基于灵敏度矩阵的方法在高导电率铁磁性 EMT 伤损检测中面临显著局限。为了摆脱对线性微扰模型的强依赖，近年来研究者逐渐将深度学习方法引入电学成像领域。端到端卷积神经网络（CNN）、编码器--解码器（Encoder--Decoder）架构以及生成对抗网络（GAN）等已在电磁层析成像等领域中展现出优于传统方法的非线性拟合能力。然而，当将这些方法直接应用于道岔伤损 EMT 检测时，仍存在以下两方面问题。

首先，经典 CNN 类方法通常要求输入和输出定义在固定分辨率的离散网格上。对于 EMT 成像任务，这意味着边界测量数据需要被重塑为伪图像格式，且重建结果的分辨率受限于网络设计阶段确定的网格尺寸。然而，EMT 系统的测量数据天然具有``多激励--多接收''的结构化特征（如本文系统中 $16 \times 15$ 的配对观测），简单将其展平为图像会丢失线圈之间的拓扑关系；同时，工程检测中往往需要在不同分辨率下评估伤损分布，固定网格输出缺乏灵活性。

其次，从数学本质上看，EMT 逆问题并非简单的``向量 $\to$ 向量''或``图像 $\to$ 图像''映射，而是更为一般的\textbf{函数到函数}（function-to-function）映射：即由边界上的观测函数（端口电压响应集合）推断内部区域的参数分布函数（电导率、磁导率或伤损指示函数）。算子学习（Operator Learning）范式正是面向此类无限维函数空间之间的映射而设计的。与标准神经网络不同，算子学习方法通过联合逼近函数表示与空间求值算子，能够在\textbf{任意空间坐标点}上进行查询输出，从而天然支持连续域上的成像重建\upcite{chen1995universal}。

在现有算子学习方法中，深度算子网络（Deep Operator Network, DeepONet）\upcite{lu2021learning} 和傅里叶神经算子（Fourier Neural Operator, FNO）\upcite{li2020fourier} 是两种代表性架构。FNO 通过在频域中进行全局卷积来逼近算子，在处理具有周期性边界和规则网格的问题时效率较高。然而，EMT 道岔伤损检测场景具有以下特殊之处：激励线圈与接收线圈的空间排布呈非均匀平面结构；缺陷引起的电磁扰动具有强局部性和方向敏感性；同时，本文不仅需要二维平面上的伤损热图，还需要在三维空间中任意非均匀采样点上重构差场分布。FNO 的频域全局处理方式对这些局部化、非均匀特征并不敏感，且其隐式网格依赖特性限制了在非规则采样点上的灵活求值能力\upcite{lu2022comprehensive}。

相比之下，DeepONet 采用的 Branch--Trunk 结构具有更好的适应性。\textbf{Branch 网络}负责编码离散观测（如多激励条件下的边界电压响应），将其映射为固定维度的隐式特征向量，天然契合 EMT 系统的多激励轮转测量模式；\textbf{Trunk 网络}则在任意连续空间坐标上进行查询求值，使输出完全摆脱固定网格的限制。这一结构特性使得 DeepONet 特别适合本文 EMT 成像任务中``结构化观测输入 + 连续空间查询输出''的核心需求。此外，Branch--Trunk 的解耦设计也为后续扩展多任务分支——如同一网络中并行实现二维热图检测与三维差场重构——提供了清晰的模块化基础。

综上所述，本文选择以 DeepONet 作为基础架构，并针对道岔伤损 EMT 检测的任务特点进行扩展，构建共享条件特征驱动的深度算子网络 SCF-DeepONet。下一章将详细阐述该网络的模型结构与训练策略。

\section{本章小结}

本章围绕 EMT 的理论基础和传统成像方法展开论述。首先介绍了 EMT 系统的基本组成与测量流程；随后从麦克斯韦方程出发，分析了 EMT 正问题的电磁学本质及其有限元求解思路；在此基础上，给出了 EMT 逆问题的离散数学表达，并从软场效应、欠定性和噪声敏感性等方面分析了其病态性来源。系统梳理了基于灵敏度矩阵的 LBP、正则化与迭代等传统成像方法，指出其在高导电率铁磁性金属伤损检测场景下的局限性。最后，从函数映射的数学本质出发，论证了算子学习范式相对于传统 CNN 和 FNO 的适配性，并阐明了选择 DeepONet Branch--Trunk 结构作为本文基础架构的理由，为第三章 SCF-DeepONet 的提出奠定了方法论基础。